%! AP Statistics — Unit 1, Lesson 7
\documentclass[10pt]{article}
\usepackage[margin=0.6in]{geometry}
\usepackage{xcolor}
\usepackage[most]{tcolorbox}
\usepackage{enumitem}
\usepackage{multicol}
\usepackage{amsmath,amssymb}
\usepackage{array}
\usepackage{tabularx}
\tcbuselibrary{breakable}

\definecolor{sky}{HTML}{EAF3FB}
\definecolor{navy}{HTML}{1F3A5F}
\definecolor{redbox}{HTML}{FCECEC}
\definecolor{greenbox}{HTML}{EDF8F0}
\definecolor{goldbox}{HTML}{FFF7E6}
\newtcolorbox{skillbox}[2][]{
    colback=#2, colframe=navy, boxrule=0.8pt, arc=2mm,
    left=2mm,right=2mm,top=1.5mm,bottom=1.5mm,
    fonttitle=\bfseries, title={#1}, halign title=center, breakable
}
\setlist[itemize]{leftmargin=1.2em,itemsep=0.35em,topsep=0.35em}
\setlist[enumerate]{leftmargin=1.4em,itemsep=0.35em,topsep=0.35em}
\newcolumntype{C}[1]{>{\centering\arraybackslash}p{#1}}
\newcolumntype{Y}{>{\centering\arraybackslash}X}

\begin{document}
\def\CourseName{AP Statistics}
\def\SchoolYear{2026--2027}
\def\MeetingLength{55 minutes}
\def\UnitNumberName{Unit 1: Exploring One-Variable Data \quad}
\def\LessonNumberName{Lesson 1.7: Summary Statistics for a Quantitative Variable}

\begin{center}
    {\Large\bfseries \CourseName: \SchoolYear\ (\MeetingLength) \\
    {\small\bfseries \UnitNumberName \LessonNumberName}}
\end{center}

\begin{tcolorbox}[colback=sky,colframe=navy,boxrule=0.9pt,arc=2mm,
    left=2mm,right=2mm,top=1.5mm,bottom=1.5mm]
    \textbf{Primary Objective:} Students will calculate and interpret the mean, median, range,
    interquartile range (IQR), and standard deviation for a quantitative variable. Students will
    select the appropriate measure of center and spread based on the shape of the distribution,
    and will identify outliers using the $1.5 \times \text{IQR}$ rule (Big Idea: UNC; AP Skills 2 \& 4).
\end{tcolorbox}

\begin{skillbox}[Priority Ideas \& Skills]{goldbox}
\begin{minipage}[t]{0.48\linewidth}
    \textbf{AP Skills 2 \& 4 --- Data Analysis \& Statistical Argumentation}
    \begin{itemize}
        \item Calculate mean, median, range, IQR, and standard deviation by hand and with technology.
        \item Choose median + IQR for skewed or outlier-affected distributions; mean + standard deviation for symmetric distributions.
        \item Apply the $1.5 \times \text{IQR}$ fence rule to identify outliers.
        \item Interpret statistics in context (e.g., ``The median commute time was 28 minutes'').
    \end{itemize}
\end{minipage}
\hfill
\begin{minipage}[t]{0.48\linewidth}
    \textbf{Key Understandings}
    \begin{itemize}
        \item Mean is pulled toward the tail in skewed distributions; median is resistant.
        \item Standard deviation measures typical distance from the mean --- it is always $\geq 0$.
        \item IQR captures the middle 50\% of data; it is resistant to outliers.
        \item The choice of summary statistics should match the shape of the distribution.
    \end{itemize}
\end{minipage}
\end{skillbox}

\begin{skillbox}[{Vocabulary, Concepts, \& Theorems}]{greenbox}
\begin{minipage}[t]{0.48\linewidth}
    \begin{itemize}
        \item \textbf{Mean} ($\bar{x}$): $\bar{x} = \dfrac{\sum x_i}{n}$; the arithmetic average; not resistant to outliers
        \item \textbf{Median} ($M$): middle value when sorted; resistant to outliers
        \item \textbf{Range}: $\max - \min$; simple but not resistant
        \item \textbf{Quartiles} ($Q_1$, $Q_2$, $Q_3$): values at the 25th, 50th, 75th percentiles
        \item \textbf{IQR}: $Q_3 - Q_1$; resistant measure of spread; contains middle 50\% of data
    \end{itemize}
\end{minipage}
\hfill
\begin{minipage}[t]{0.48\linewidth}
    \begin{itemize}
        \item \textbf{Standard deviation} ($s_x$): $s_x = \sqrt{\dfrac{\sum(x_i - \bar{x})^2}{n-1}}$; typical distance from the mean; not resistant
        \item \textbf{Variance} ($s_x^2$): square of the standard deviation
        \item \textbf{Outlier rule}: a value is an outlier if it is below $Q_1 - 1.5\,\text{IQR}$ or above $Q_3 + 1.5\,\text{IQR}$
        \item \textbf{Resistant statistic}: not strongly affected by extreme values (median, IQR)
    \end{itemize}
\end{minipage}
\end{skillbox}

\begin{skillbox}[Activate Prior Knowledge \& Spiral Review]{sky}
\begin{minipage}[t]{0.48\linewidth}
    \textbf{Quick Review (3 min)}\\
    Revisit the SOCS framework from Lesson 1.6. Ask: ``When I said `center is approximately 8 hours,' how could I be more precise?'' Students recall mean and median from prior coursework. Today we formalize and extend.
\end{minipage}
\hfill
\begin{minipage}[t]{0.48\linewidth}
    \textbf{Spiral Connections}
    \begin{itemize}
        \item Lesson 1.8: five-number summary ($\min, Q_1, M, Q_3, \max$) is displayed in a boxplot.
        \item Lesson 1.9: comparing distributions uses mean/median and standard deviation/IQR.
        \item Unit 5--9: standard deviation plays a central role in all inference procedures.
    \end{itemize}
\end{minipage}
\end{skillbox}

\begin{skillbox}[Hook \& Lesson]{sky}
\begin{minipage}[t]{0.48\linewidth}
    \textbf{Hook (5 min)}\\
    Scenario: 5 friends' salaries are \$30k, \$32k, \$35k, \$38k, \$200k. Ask: ``What is the typical salary?'' Compute the mean (\$67k) and the median (\$35k). Discuss: ``Which is more representative?'' --- introduces the resistance concept immediately.
\end{minipage}
\hfill
\begin{minipage}[t]{0.48\linewidth}
    \textbf{Lesson Flow (18 min)}
    \begin{enumerate}
        \item Calculate mean and median for two datasets: one symmetric, one skewed.
        \item Show algebraically why skewness pulls mean toward the tail.
        \item Compute $Q_1$, $Q_3$, IQR; apply the outlier fence rule.
        \item Calculate standard deviation step-by-step for a small dataset; connect to the idea of ``average spread.''
        \item Summarize the decision rule: skewed/outliers $\to$ median + IQR; symmetric $\to$ mean + $s_x$.
    \end{enumerate}
\end{minipage}
\end{skillbox}

\begin{skillbox}[Explicit Instruction \& Active Monitoring]{sky}
\begin{minipage}[t]{0.48\linewidth}
    \textbf{Direct Instruction (10 min)}
    \begin{itemize}
        \item Walk through the standard deviation formula slowly; emphasize: squaring deviations makes all positive, square root returns units.
        \item Demonstrate with technology (calculator or software) after the by-hand calculation.
        \item Outlier fence: lower = $Q_1 - 1.5(Q_3-Q_1)$; upper = $Q_3 + 1.5(Q_3-Q_1)$.
    \end{itemize}
\end{minipage}
\hfill
\begin{minipage}[t]{0.48\linewidth}
    \textbf{Active Monitoring}
    \begin{itemize}
        \item Check: Are students using $n-1$ (sample standard deviation $s_x$) rather than $n$?
        \item Common error: confusing the median finding rule for even vs.\ odd $n$.
        \item Ask: ``If I add a very large value to this dataset, which statistic changes more --- the mean or the median? The IQR or the standard deviation?''
    \end{itemize}
\end{minipage}
\end{skillbox}

\begin{skillbox}[Group Work \& Differentiation]{redbox}
\begin{minipage}[t]{0.48\linewidth}
    \textbf{Group Activity (8--10 min)}\\
    Groups receive a dataset (e.g., 12 recorded lap times). Each group:
    \begin{enumerate}
        \item Calculates mean, median, range, IQR.
        \item Checks for outliers using the fence rule.
        \item Determines which pair (mean+SD or median+IQR) is more appropriate.
        \item Writes one sentence choosing and justifying their statistic.
    \end{enumerate}
\end{minipage}
\hfill
\begin{minipage}[t]{0.48\linewidth}
    \textbf{Differentiation}
    \begin{itemize}
        \item \textit{Support}: Provide a step-by-step computation sheet for standard deviation; allow calculator use for all calculations after doing one by hand.
        \item \textit{Extension}: Explore how adding a constant to every value, or multiplying every value by a constant, affects each statistic. (Leads to linear transformation rules in Lesson 1.10.)
        \item \textit{Connections}: Discuss why insurance companies, economists, and scientists might prefer different statistics.
    \end{itemize}
\end{minipage}
\end{skillbox}

\begin{skillbox}[Individual Work \& Assessment]{redbox}
\begin{minipage}[t]{0.48\linewidth}
    \textbf{Exit Ticket (5 min)}\\
    Dataset: 4, 7, 7, 9, 10, 13, 15, 30. Find: (1) mean, (2) median, (3) IQR, (4) check if 30 is an outlier using the fence rule, (5) state which measure of center is more appropriate and why.
\end{minipage}
\hfill
\begin{minipage}[t]{0.48\linewidth}
    \textbf{AP-Style Check}\\
    \textit{``The mean and median of a distribution are 45 and 32, respectively. What does this suggest about the shape of the distribution? Explain.''}\\[4pt]
    Assess: student connects mean $>$ median $\to$ skewed right, with explanation.
\end{minipage}
\end{skillbox}

\begin{skillbox}[Reinforcement \& Extension]{goldbox}
\begin{minipage}[t]{0.48\linewidth}
    \textbf{Homework / Practice}
    \begin{itemize}
        \item Compute all five summary statistics for two datasets; identify outliers; state appropriate pair of statistics for each.
        \item AP textbook: multiple-choice and free-response on summary statistics.
    \end{itemize}
\end{minipage}
\hfill
\begin{minipage}[t]{0.48\linewidth}
    \textbf{Extension / Preview}
    \begin{itemize}
        \item \textit{Extension}: Investigate how the standard deviation changes when one outlier is added or removed from a dataset.
        \item \textit{Preview}: Lesson 1.8 --- the five-number summary ($\min, Q_1, M, Q_3, \max$) is visualized in a \textit{boxplot}. Think about what information a boxplot would show that a histogram does not.
    \end{itemize}
\end{minipage}
\end{skillbox}

\end{document}
